
% --------------------------------------------------------------------------------------------------
% Kommando Beispiele
% --------------------------------------------------------------------------------------------------

% Referenz auf z.B. Kapitel
\chapter{Foundations of Proof Assistants}
\label{c:chapter3}

\section{Systems of Logic}

In classical logic, statements (=Propositions) are either true or false. No test or proof is needed whether a proposition is true or false. E.g. having a proposition $A$ the disjunction $A \lor \neg A$ is always true without showing which side of the disjunction holds. The Law of the Excluded Middle (LEM) holds. \\

On the other side, intuitionistic logic
rejects the LEM. Claiming $A \lor \neg A$ requires to prove at least one side of the disjunction namely $A$ or $\neg A$.\\
Proving $\exists \alpha : P(\alpha)$ requires a witness for $\alpha$ and an algorithm (or program) to compute $P(\alpha)$. A theorem in intutioniscit logic implies beeing a theorem in classical logic but not in the opposite directo. in general a tautology in classical logic is not a theorem in intutionistic logic.\\

\textit{Question: Lean 4 seems to base on Intuistic logic and has Proof by contradiction, which is not a constructive proof becaust is can be usd to proof existencewithutproviding an exampleof the obj that is proves existe. \cite{web:intuitionistic_logic}}

Why this matters? Intuitionistic logic requires a constructive proof that an object exists and this construction suits sas an anlgorithm or programm to generate such object.the  Curry-howjard correspondence defines an isomorphism between such programms and proofs. As a consequence, a program in lean 4 or other proof assistants can be used as a proof and vice versa. \cite{web:intuitionistic_logic}


\section{Type Theory}
... type theory and typed $\lambda$ - calculus shall be before the Curry-Howard Correspondence

\section{Curry–Howard Correspondence}

From Wikipedia\\

$$ formula	\iff type $$
$$proof \iff 	term $$
$$ \text{formula is true} \iff	\text{type has an element} $$

from https://dev.to/shrsv/from-logic-to-code-understanding-the-curry-howard-correspondence-in-lean-1and \\
the CHc reveals that
\begin{itemize}
    \item logical propositions $\iff$ types
    \item proofs $\iff$ programs 
    \item proof checking $\iff$ type checking
\end{itemize}
As type checking is done by a compiler, a such compler is capable to 
perform proof checking. \\
Lets look at an example. $\P \land Q \to P$ tells: if both $P$ and $Q$ 
are true, then $P$ is true. To do so, one must prove $P$ and $Q$, or 
similarly provide a proof for both. To prove $P \to Q$ one must
provide
\begin{itemize}
    \item a mapping that transfroms a proof of $P$ to a proof of $Q$
\end{itemize}
and this is what functions do. 


\section{Why should we thrust Lean 4}
To trust such a mechanical verification process, it is enough to show that the verification algorithm actually verifies that all formal rules are correctly applied. \cite[Chapter~1.2]{coq_book}
